\begin{abstract}
    Predicting the directional movement of financial assets remains a foundational challenge in quantitative finance due to the non-stationarity and high-entropy nature of market data. This research presents a \textbf{Multi-Source Financial Intelligence System} that integrates heterogeneous data streams---including historical price action, macroeconomic indicators, and transformer-based sentiment analysis---to mitigate information asymmetry. Utilizing \textbf{FinBERT}, we extract daily sentiment polarity from news headlines to serve as a proxy for behavioral market drivers. The modeling pipeline employs an \textbf{XGBoost} ensemble, scientifically validated through Augmented Dickey-Fuller (ADF) tests and rigorous Out-of-Sample (OOS) backtesting. Crucially, our evaluation framework accounts for **market friction (10 bps transaction costs)** to ensure results represent executable performance. Our findings reveal that while most high-liquidity assets closely follow a random-walk behavior under friction, the fusion of behavioral sentiment and technical indicators identifies robust momentum in commodities (Gold), achieving an OOS Sharpe Ratio of 1.49. This study provides a scientific framework for trust-based machine learning, emphasizing that true predictive alpha must be measured against realistic execution constraints.
\end{abstract}